% !PRINTED 440
\weberpage{440}
%<<CONT
\gesperrt{The totality of the numbers $(\varrho)$ thus agrees, apart from sign and from order, with the totality of the numbers $(\nu)$.}

In the formulae (1) and (2) we can, however, on account of the periodicity of sine and cosine, replace $\nu m$ by $\varrho$, and since $\cos(-\varphi) = \cos\varphi$, we can in formula (1) also replace $m\nu$ by $\nu$, i.e. the sign does not depend on $m$. To determine it, we take into account that the cosine of an angle in the first quadrant is positive, in the second quadrant negative, and thus that the sign in (1) is positive or negative according as an even or an odd number of the angles
\[
\frac{2\nu\pi}{n}
\]
lies in the second quadrant, or according as an even or an odd number of the numbers $\nu$ lies between $\dfrac{n}{4}$ and $\dfrac{n}{2}$.

If $x$ is any number that is not an integer, we denote by $E(x)$ the greatest integer that is less than $x$, so that $x$ lies between $E(x)$ and $E(x)+1$. The number of integers lying between two numbers $x$ and $y$ is then equal to $E(y) - E(x)$, and we therefore have to investigate whether
\[
E\left(\frac{n}{2}\right) - E\left(\frac{n}{4}\right)
\]
is an even or an odd number. We must distinguish four cases, as appears from the following compilation, in which $k$ denotes a non-negative integer:
\[
\begin{array}{lll}
n = 8k+1, & E\left(\dfrac{n}{2}\right) = 4k, & E\left(\dfrac{n}{4}\right) = 2k, \\[2ex]
n = 8k+3, & E\left(\dfrac{n}{2}\right) = 4k+1, & E\left(\dfrac{n}{4}\right) = 2k, \\[2ex]
n = 8k+5, & E\left(\dfrac{n}{2}\right) = 4k+2, & E\left(\dfrac{n}{4}\right) = 2k+1, \\[2ex]
n = 8k+7, & E\left(\dfrac{n}{2}\right) = 4k+3, & E\left(\dfrac{n}{4}\right) = 2k+1.
\end{array}
\]
Thus in (1) the positive sign is to be taken when $n$ is of the form $8k+1$ or $8k+7$, the negative when $n$ is of the form $8k+3$ or $8k+5$. In the first cases
%CONT>>
